\chapter{RESULTADOS EXPERIMENTALES}
\label{ch:resultados}

Este capítulo presenta los resultados obtenidos durante la validación del prototipo de robot móvil diferencial. La odometría diferencial y la sintonización PID se evaluaron en simulación y en hardware real; el mapeo con \texttt{slam\_toolbox}, la navegación autónoma y la exploración de fronteras se evaluaron sobre el hardware real, en el entorno físico disponible. Los resultados complementarios de exploración y mapeo se presentan en el \ref{anx:protocolo_resultados}. El entorno físico en que se realizaron las pruebas sobre hardware real se muestra en la Fig.~\ref{fig:entorno_experimental}.

\utemFigurePropia{2.Texto/Imag/entorno_experimental.jpg}{0.78\textwidth}
{Entorno utilizado para las pruebas experimentales en hardware real.}{fig:entorno_experimental}

\section{Evaluación de odometría y control diferencial}
\label{sec:resultados_odometria}

La evaluación rectilínea utiliza una referencia nominal de 2~m. El script publica una velocidad lineal de 0.13~m/s durante $2/0.13\approx15.4$~s y luego detiene el robot, por lo que la referencia no es una consigna de posición sino el producto de la velocidad comandada por la duración del comando. En hardware real, la distancia física se midió con cinta métrica y el error odométrico corresponde a la discrepancia entre la odometría de encoders y esa medición. En simulación, la consulta de la pose verdadera a Gazebo no entregó datos durante los ensayos y el script registró la propia odometría como distancia de respaldo. Por ello, en simulación no se dispone de una verdad de terreno independiente.

Las pruebas angulares se analizaron respecto de una referencia nominal de
$360^\circ$ definida por el protocolo, que no corresponde a una consigna
efectiva de vuelta completa; por ello, sus resultados expresan desviación
respecto de ese valor nominal y no un error odométrico angular aislado. El
desarrollo del protocolo y el límite que origina esta distinción se conservan
en el \ref{anx:odometria_detalle}.
Para distinguir ambos fenómenos se utilizaron las siguientes métricas:

\begin{align}
E_{\mathrm{seg}}
&=
\frac{
\left|d_{\mathrm{real}}-d_{\mathrm{ref}}\right|
}{
d_{\mathrm{ref}}
}\,100\%,
\label{eq:error_seguimiento}\\
E_{\mathrm{odom}}
&=
\frac{
\left|d_{\mathrm{odom}}-d_{\mathrm{real}}\right|
}{
d_{\mathrm{real}}
}\,100\%.
\label{eq:error_odometrico}
\end{align}

\noindent donde $d_{\mathrm{ref}}$ es la distancia comandada (m), $d_{\mathrm{real}}$ la distancia física medida o verdad de terreno (m),
$d_{\mathrm{odom}}$ la distancia calculada mediante odometría (m),
$E_{\mathrm{seg}}$ el error de seguimiento (\%) y
$E_{\mathrm{odom}}$ el error odométrico (\%).
La desviación angular normalizada se expresa respecto de esa referencia nominal fija:

\begin{equation}
E_{\theta,360}
=
\frac{
\left|
\Delta\theta_{\mathrm{obs}}
-
\Delta\theta_{360}
\right|
}{
\left|\Delta\theta_{360}\right|
}
\,100\%,
\qquad
\Delta\theta_{360}=360^\circ.
\label{eq:desviacion_angular_360}
\end{equation}

\noindent donde $E_{\theta,360}$ es la desviación angular normalizada [\%],
$\Delta\theta_{360}$ la referencia nominal fija y
$\Delta\theta_{\mathrm{obs}}$ el giro observado. La
Tabla~\ref{tab:resumen_odometria} consolida las métricas diferenciadas
obtenidas en la trayectoria rectilínea de 2~m.

\begin{table}[H]
\centering
\caption{Métricas diferenciadas de seguimiento y odometría en la trayectoria rectilínea de 2~m.}
\label{tab:resumen_odometria}
\small
\begin{tabular}{p{7.2cm} p{2.2cm} p{2.8cm}}
\toprule
\textbf{Métrica} & \textbf{Simulación} & \textbf{Hardware real} \\
\midrule
Error medio de seguimiento $E_{\mathrm{seg}}$
& 3.39\%$^{*}$ & 7.18\% \\
Desviación estándar de $E_{\mathrm{seg}}$ ($n=5$)
& 2.04\%$^{*}$ & 0.43\% \\
Error medio odométrico $E_{\mathrm{odom}}$
& NA & 3.63\% \\
Desviación estándar de $E_{\mathrm{odom}}$ ($n=5$)
& NA & 0.44\% \\
\bottomrule
\end{tabular}
\vspace{3pt}
\parbox{0.85\linewidth}{\footnotesize $^{*}$ En simulación se calcula con la distancia de odometría, porque la pose verdadera de Gazebo no quedó registrada.}
\end{table}

En hardware real, el error medio de seguimiento de 7.18\% corresponde a una
distancia física media de aproximadamente 1.856~m frente a los 2~m definidos
por el protocolo. Este error debe distinguirse del error odométrico de 3.63\%:
el primero caracteriza cuánto se apartó el desplazamiento físico de la
referencia comandada, mientras que el segundo compara la estimación de
odometría con la distancia realmente recorrida.

Los registros permiten separar el 7.18\% en dos partes, aunque no aislar sus
causas físicas. La odometría media de las cinco repeticiones fue 1.924~m, es
decir, los encoders registraron un avance 3.8\% menor que los 2~m nominales.
Esta parte ya está presente en la medición de los encoders y es compatible con
el tiempo de subida de las ruedas (Tabla~\ref{tab:pid_performance}), porque la
referencia se define por tiempo y no por posición. La simulación, que no
reproduce la zona muerta del driver, presenta una desviación del mismo orden
(3.39\%). La parte restante, 0.067~m o 3.4\% de la referencia, corresponde a la
sobreestimación de la distancia por la odometría, que expresada respecto de la
distancia real es el error odométrico de 3.63\%. Su baja dispersión (0.44\%)
indica un sesgo sistemático, compatible con un error de escala en el radio
efectivo de rueda o con deslizamiento persistente. Estos efectos no se
separaron mediante ensayos independientes. El valor $N_e=1980$ no explica el
sesgo, porque su diferencia con los promedios medidos haría que la odometría
subestimara la distancia entre 0.2\% y 0.6\%.

Los resultados correspondientes al protocolo angular se resumen en la Tabla~\ref{tab:resumen_giro}; su comparación entre simulación y hardware se desarrolla en la Sec.~\ref{sec:comparacion_sim_real}.

\begin{table}[H]
\centering
\caption{Desviación angular y error medio respecto de la referencia nominal de $360^\circ$ (5 repeticiones por entorno).}
\label{tab:resumen_giro}
\small
\begin{tabular}{p{7.0cm} p{3.2cm} p{3.2cm}}
\toprule
\textbf{Métrica} & \textbf{Simulación} & \textbf{Hardware real} \\
\midrule
Desviación angular media & $17.1^\circ$ & $11.7^\circ$ \\
Desviación angular normalizada $E_{\theta,360}$ & $4.74\%$ & $3.24\%$ \\
Desviación estándar & $5.8^\circ$ ($1.61\%$) & $8.0^\circ$ ($2.22\%$) \\
\bottomrule
\end{tabular}
\end{table}

Esta métrica no es un error odométrico angular aislado. En simulación, el giro observado es el giro acumulado de \texttt{/odom} al final de un ensayo de duración fija. En hardware real, es el giro físico medido al final del ensayo. En ambos casos se compara con la referencia nominal de 360~grados y no con el giro efectivamente comandado. Los giros observados, entre $335.9^\circ$ y $347.4^\circ$ en simulación y entre $340^\circ$ y $380^\circ$ en hardware real, superan los $315^\circ$ que se obtendrían si el límite de 0.35~rad/s actuara de forma ideal y constante durante todo el ensayo (\ref{anx:odometria_detalle}). Por tanto, los registros no son coherentes con esa hipótesis y la causa de la diferencia no se determinó. Un giro comandado a un ángulo fijo y controlado, planteado como trabajo futuro, entregaría un error angular directo.
El desglose por repetición de ambos protocolos (trayectoria rectilínea y
protocolo angular, en simulación y hardware real) se presenta en la
Fig.~\ref{fig:resultados_odometria_repeticiones}.

\begin{utemMultiFigurePropia}{Resultados por repetición de las pruebas rectilíneas y angulares en simulación y hardware real.}{fig:resultados_odometria_repeticiones}
\addutem{0.48}{2.Texto/Imag/odom_error_barras_sim.png}{Trayectoria rectilínea en simulación.}
\hfill
\addutem{0.48}{2.Texto/Imag/odom_error_barras.png}{Trayectoria rectilínea en hardware real.}
\\[0.4cm]
\addutem{0.48}{2.Texto/Imag/giro_error_barras_sim.png}{Protocolo angular en simulación.}
\hfill
\addutem{0.48}{2.Texto/Imag/giro_deriva_barras.png}{Protocolo angular en hardware real.}
\end{utemMultiFigurePropia}

La desviación angular observada integra los efectos del límite de
velocidad, la duración del ensayo, la dinámica física y el procedimiento de
medición. La incorporación de una IMU se plantea como mejora futura.

\section{Evaluación del controlador PID}
\label{sec:resultados_pid}

El seguimiento en lazo cerrado se verificó en dos momentos del procedimiento
de sintonización descrito en la Sec.~\ref{sec:sintonizacion_pso}: con las
ganancias sugeridas por PSO y con las ganancias definitivas obtenidas tras el
refinamiento. Ambas capturas se realizaron sobre el robot real, con una
referencia de 0.13~m/s y el mismo protocolo, conservado en el
\ref{anx:captura_lazo_cerrado}.

\utemFigurePropia{2.Texto/Imag/respuesta_escalon_final_pso.png}{0.82\textwidth}
{Respuesta en lazo cerrado del robot real con las ganancias sugeridas por PSO,
antes del refinamiento experimental (referencia de 0.13~m/s).}
{fig:respuesta_escalon_pso}

Con las ganancias del PSO, ninguna de las dos ruedas sostiene la referencia
durante todo el ensayo. La rueda izquierda se estabiliza en 0.12~m/s y solo
alcanza los 0.13~m/s de forma intermitente, manteniendo un error estacionario
cercano a 0.01~m/s. La rueda derecha sube por tramos, con una detención en
torno a 0.11~m/s antes de continuar, y recién se mantiene en la referencia
después de los 6~s. En ambos casos la subida es más lenta que la obtenida con
las ganancias definitivas.

\utemFigurePropia{2.Texto/Imag/respuesta_escalon_final.png}{0.82\textwidth}
{Respuesta experimental en lazo cerrado del robot real con las ganancias
definitivas y referencia de 0.13~m/s.}
{fig:respuesta_escalon_final}

Con las ganancias definitivas desaparece el error estacionario de la rueda
izquierda y la detención intermedia de la derecha, y la referencia se alcanza
cerca de un segundo antes. En ambas ruedas, la velocidad medida converge hacia la referencia de
0.13~m/s. La rueda derecha presenta una respuesta más uniforme, mientras que
la izquierda muestra pequeñas variaciones discretas alrededor del valor
objetivo. La figura constituye una verificación cualitativa del seguimiento
sobre el robot real. Los indicadores de la Tabla~\ref{tab:pid_performance}
se calcularon sobre la captura de la Fig.~\ref{fig:respuesta_escalon_final}, con la referencia de 0.13~m/s. La
comparación ilustrativa con el modelo FOPDT y la sintonía Lambda se conserva
en el \ref{anx:lambda_detalle}. En esa simulación, las ganancias finales producen sobreimpulsos de 10.7\% a 15.6\%. Esos valores provienen de un modelo que no reproduce la zona muerta del driver, las bandas muertas asimétricas ni los tipos enteros del firmware, por lo que no se comparan directamente con la Tabla~\ref{tab:pid_performance}.

\begin{table}[H]
\centering
\caption{Indicadores de la respuesta al escalón en lazo cerrado, medidos sobre la captura experimental del robot real con referencia de 0.13~m/s.}
\label{tab:pid_performance}
\small
\begin{tabular}{p{6.8cm} p{3.1cm} p{3.1cm}}
\toprule
\textbf{Indicador} & \textbf{Rueda izquierda} & \textbf{Rueda derecha} \\
\midrule
Velocidad estacionaria medida ($t\geq 7$~s) & 0.128 m/s & 0.130 m/s \\
Error estacionario & 0.002 m/s & 0.000 m/s \\
Tiempo de subida 10\%--90\% & 0.8 s & 1.4 s \\
Sobreimpulso máximo & Sin sobreimpulso significativo & Sin sobreimpulso significativo \\
\bottomrule
\end{tabular}
\end{table}

Ambas ruedas alcanzan la referencia sin sobreimpulso apreciable. La velocidad informada por el firmware tiene una resolución de 0.01~m/s, cercana a un 8 por ciento de la referencia. Por eso la banda del 2 por ciento queda por debajo de esa resolución y no se reporta un tiempo de establecimiento al 2 por ciento. La respuesta se estabiliza en torno a la referencia dentro de los primeros segundos, dentro de la resolución de medida.
La comparación completa entre las ganancias del firmware y la referencia
analítica Lambda se conserva en la Fig.~\ref{fig:anx_comparacion_pso_vs_lambda}, con carácter ilustrativo.

La ganancia integral difiere entre ruedas, con $K_{i,L}=0.075$ y $K_{i,R}=0.001$. El valor izquierdo se fijó durante el refinamiento experimental, porque con las ganancias sugeridas por PSO esa rueda mantenía un error estacionario bajo la referencia (Fig.~\ref{fig:respuesta_escalon_pso}). El efecto de estas ganancias está condicionado por el almacenamiento entero de \texttt{ITerm} (\ref{anx:pid_ecuaciones}). En la rueda derecha, el producto $K_{i,R}e_R^{\mathrm{db}}[k]$ solo alcanzaría una unidad con errores de 1000~ticks/frame, muy superiores a los del ensayo, por lo que su estado integral permanece en cero y $K_{i,R}$ no tiene efecto en la ley ejecutada. En la rueda izquierda, el estado integral solo puede crecer en magnitud en ciclos con $|e_L^{\mathrm{db}}[k]|\geq14$~ticks/frame, cerca de 0.045~m/s. En el escalón de 0.13~m/s, de unos 41~ticks/frame, esto ocurre principalmente durante el arranque, y el valor acumulado queda sumado al numerador del incremento de control. Por ello, la mejora observada con las ganancias definitivas se informa como un resultado empírico del conjunto de ganancias y no como el efecto de una acción integral convencional.

El PSO entregó un punto de partida heurístico para las tres ganancias. El refinamiento experimental ajustó ese punto de partida, con aumentos de $K_p$ de 28\% y 63\%, reducciones de $K_d$ de 41\% y 45\% y un aumento de varios órdenes de magnitud en $K_i$, hasta lograr el seguimiento de la referencia sobre el robot real. Por tanto, el aporte del PSO fue el de una semilla de búsqueda y no el de un ajuste final validado por sí solo.
El retardo identificado resultó menor que el período de captura, por lo que no introduce un desplazamiento temporal efectivo en el modelo discreto. En la práctica el modelo se comporta como uno de primer orden. El detalle se conserva en el \ref{anx:fodt_scripts}.

\section{Resultados de mapeo con \texttt{slam\_toolbox}}
\label{sec:resultados_slam}

El mapeo se evaluó sobre el hardware real a lo largo de varias sesiones prácticas dentro del entorno experimental, de las cuales cinco iteraciones produjeron un mapa válido. La inspección visual de los mapas obtenidos mostró paredes definidas y consistencia global. Esta valoración es cualitativa. El análisis \emph{post hoc} no incorpora una detección automática de cierre de lazo, por lo que ese indicador quedó registrado como no determinado en las cinco iteraciones válidas.
La métrica utilizada corresponde al índice de cobertura libre, definido como
la proporción de celdas observadas como libres respecto del total de celdas
susceptibles de ser exploradas:

\begin{equation}
C_{\mathrm{libre}}
=
100\,
\frac{N_{\mathrm{libre}}}
{N_{\mathrm{libre}}+N_{\mathrm{desconocido}}},
\label{eq:cobertura_libre}
\end{equation}

\noindent donde $N_{\mathrm{libre}}$ es el número de celdas clasificadas como
libres, $N_{\mathrm{desconocido}}$ el número de celdas desconocidas y
$C_{\mathrm{libre}}$ el índice de cobertura libre (\%). Las celdas ocupadas no
se incluyen en el denominador. La clasificación se realiza sobre la imagen de
ocupación generada por \texttt{slam\_toolbox}, según los umbrales de
intensidad indicados en el \ref{anx:resultados_metodologia}.
El área libre mapeada se calcula a partir del número de celdas libres:

\begin{equation}
A_{\mathrm{libre}}
=
N_{\mathrm{libre}}\,\rho_{\mathrm{map}}^2,
\label{eq:area_libre_mapeada}
\end{equation}

\noindent donde $A_{\mathrm{libre}}$ es el área libre mapeada [m$^2$] y
$\rho_{\mathrm{map}}$ la resolución del mapa [m/celda]. La
Tabla~\ref{tab:slam_real_indicadores} consolida los indicadores de mapeo
obtenidos en las cinco iteraciones válidas sobre hardware real, recuperados \emph{post hoc} del \textit{rosbag} y del mapa PGM de cada sesión; su procedencia y
alcance se documentan en el \ref{anx:resultados_metodologia}.

\begin{table}[H]
\centering
\caption{Indicadores de mapeo en hardware real por iteración (resolución $0.05$~m/celda).}
\label{tab:slam_real_indicadores}
\begin{threeparttable}
\small
\setlength{\tabcolsep}{5pt}
\renewcommand{\arraystretch}{1.15}
\begin{tabular}{
>{\centering\arraybackslash}p{1.5cm}
>{\centering\arraybackslash}p{2.5cm}
>{\centering\arraybackslash}p{2.8cm}
>{\centering\arraybackslash}p{2.8cm}
>{\centering\arraybackslash}p{2.8cm}
}
\toprule
\textbf{Iter.} & \makecell{\textbf{Dur.}\\\textbf{(min)}} & \makecell{\textbf{Dist.}\\\textbf{odom. (m)}} & \makecell{\textbf{Cobertura}\\\textbf{(\%)}} & \makecell{\textbf{Área libre}\\\textbf{(m$^2$)}} \\
\midrule
1 & 3.45 & 6.305 & 82.9 & 3.510 \\
2 & \multicolumn{4}{c}{Incompleta (sin mapa)} \\
3 & 2.16 & 3.536 & 53.5 & 3.540 \\
4 & 2.27 & 4.144 & 54.2 & 3.478 \\
5 & 2.50 & 3.850 & 67.1 & 3.483 \\
6 & 2.30 & 4.066 & 53.4 & 3.478 \\
\midrule
$\mu$ & 2.54 & 4.38 & 62.22 & 3.50 \\
$\sigma$ & 0.53 & 1.10 & 12.94 & 0.03 \\
\bottomrule
\end{tabular}
\begin{tablenotes}[flushleft]
\footnotesize
\item[] $\mu$: media aritmética. \quad $\sigma$: desviación estándar (estimador muestral, $n-1$). Ambas sobre las cinco iteraciones válidas.
\end{tablenotes}
\end{threeparttable}
\end{table}

Estos indicadores caracterizan el mapeo obtenido en las cinco iteraciones
válidas de exploración y no constituyen una medida del error absoluto del mapa. El contraste cualitativo de los mapas obtenidos se documenta en la Tabla~\ref{tab:diferencias_mapas} del \ref{anx:resultados_metodologia}.
El índice de cobertura libre se calcula sobre la imagen de ocupación completa, cuyos límites no están anclados al recinto físico. Por eso es un indicador interno relativo y no una fracción absoluta del área del recinto. Esto explica que el área libre mapeada se mantenga casi constante entre iteraciones mientras el índice de cobertura varía de forma amplia, ya que este último depende de cuántas celdas desconocidas quedan dentro de la imagen. Por su mayor estabilidad, el área libre mapeada y el resultado de las metas de la Tabla~\ref{tab:anx_explore_metas} se toman como evidencia principal de la exploración, y el índice de cobertura como indicador complementario.

\section{Navegación autónoma y exploración}
\label{sec:resultados_nav2}

La navegación dirigida se evalúa funcionalmente, mientras que la exploración
se analiza mediante cobertura y área mapeada.
La exploración por fronteras se ejecutó de forma autónoma en cinco iteraciones válidas. El nodo explorador seleccionó fronteras y despachó metas de navegación a Nav2 mediante la acción \texttt{NavigateToPose}. Los contadores de texto de \texttt{run\_explore.sh} no distinguen metas únicas, por eso informaban cero metas generadas. Para obtener una medida fiable, cada meta se reconstruyó por su identificador único a partir del flujo de estado de la acción registrado en cada sesión, según la Tabla~\ref{tab:anx_explore_metas} del \ref{anx:protocolo_resultados}.

En total se generaron 60 metas de navegación. De ellas, 59 alcanzaron el estado de ejecución, es decir el robot navegó activamente hacia la frontera. El resultado predominante fue el aborto, con 54 metas abortadas y una sola meta completada con éxito. Las metas abortadas tuvieron una vida mediana de entre 20 y 61 segundos según la iteración, de modo que corresponden a intentos reales de navegación y no a reemplazos inmediatos de meta. Otras cinco metas quedaron sin estado terminal al finalizar sus respectivos registros. Por tanto, la fracción de metas exitosas respecto del total generado fue de $1/60\approx1.7\%$, mientras que entre las 55 metas con resultado terminal fue de $1/55\approx1.8\%$. Estos valores muestran que la cadena explorador--Nav2 se ejecutó, pero también evidencian un desempeño operacional limitado para completar las metas seleccionadas en el entorno utilizado.

Los diecisiete mensajes de lista de fronteras agotada registrados en las cinco
iteraciones válidas, once de ellos en la iteración~1, no representan diecisiete
abortos ni deben compararse directamente con las 54 metas abortadas. En el código
ejecutado, una frontera puede incorporarse a la lista negra cuando una acción
de navegación retorna \texttt{ABORTED} o cuando se supera
\texttt{progress\_timeout}. El mensaje de lista agotada aparece únicamente
cuando todas las fronteras candidatas del instante coinciden con elementos de
la lista negra, momento en que esta se vacía para permitir nuevos intentos.
Por ello, ambos contadores describen eventos diferentes.
En conjunto, los ensayos verifican la integración funcional del proceso de
exploración: se detectaron fronteras, se generaron metas y Nav2 intentó
ejecutarlas mientras el mapa continuó actualizándose. Sin embargo, la elevada
proporción de metas abortadas impide interpretar estos ensayos como evidencia
de navegación robusta hacia las fronteras. El resultado se conserva como una
limitación cuantificada del desempeño obtenido en el entorno experimental.

\begin{utemMultiFigurePropia}{Ejecución de navegación autónoma y exploración por fronteras sobre hardware real: avance desde un mapa parcial hasta una representación de mayor cobertura, con la ruta planificada por Nav2 (trazo verde).}{fig:nav_exploracion_real}
\addutem{0.46}{2.Texto/Imag/prueba_real_navegacion_y_exploracion_parte1.png}{Mapa parcial en una etapa temprana de la exploración.}
\hfill
\addutem{0.46}{2.Texto/Imag/prueba_real_navegacion_y_exploracion_parte4.png}{Mayor cobertura con ruta planificada hacia la frontera seleccionada.}
\end{utemMultiFigurePropia}

La mejor iteración alcanzó un índice de cobertura libre de 82.9\%. El
radio de inflación de 0.15~m aplicado en un recinto reducido
(Sec.~\ref{sec:costmaps_ch9}) condiciona la maniobrabilidad. En conjunto, los ensayos verifican la integración funcional, cuantifican la
cobertura y documentan el resultado de las metas de exploración según la
Tabla~\ref{tab:anx_explore_metas}. La finalización correcta de las cinco
iteraciones se refiere a la generación de un mapa válido. La navegación dirigida con metas fijas definidas por el operador, distinta de
las metas que genera el explorador, queda fuera del alcance declarado en la
\hyperref[sec:alcances]{Alcances y limitaciones} y no se ensayó, por lo que este trabajo no informa
tasa de éxito ni error de posición final para ese caso. La Tabla~\ref{tab:estado_validacion} sintetiza el estado de validación alcanzado por cada subsistema.

\begin{table}[H]
\centering
\caption{Estado de validación por subsistema.}
\label{tab:estado_validacion}
\small
\begin{tabular}{p{2.1cm} p{5.4cm} p{5.9cm}}
\toprule
\textbf{Subsistema} & \textbf{Estado} & \textbf{Observación} \\
\midrule
Odometría & Validada cuantitativamente en hardware real (recta de 2~m, $n=5$) & En simulación no se registró verdad de terreno independiente. El protocolo angular no mide error odométrico angular \\
Control PID & Evaluado experimentalmente & Indicadores de la Tabla~\ref{tab:pid_performance} calculados sobre la captura real en lazo cerrado a 0.13~m/s \\
SLAM & Verificado funcionalmente y caracterizado cuantitativamente &
Cobertura y área libre cuantificadas en cinco iteraciones válidas; la
alineación de paredes se evaluó visualmente y no se dispuso de ground truth
geométrico para validar el error absoluto del mapa \\
Nav2 & Verificada funcionalmente, resultado de metas cuantificado & Distribución de resultados de 55 metas terminales de exploración (Tabla~\ref{tab:anx_explore_metas}). El protocolo con metas fijas queda fuera del alcance declarado (\hyperref[sec:alcances]{Alcances y limitaciones}) \\
Exploración & Verificada funcionalmente, resultado de metas cuantificado & 60 metas autónomas en 5 iteraciones, 54 abortadas, 1 exitosa y 5 sin estado terminal. Cobertura y área cuantificadas \\
Localización AMCL & Implementada y configurada, no evaluada experimentalmente &
El nodo y sus archivos de lanzamiento están disponibles para navegación sobre
mapas almacenados, pero AMCL no se ejecutó durante los ensayos de mapeo y
exploración ni se realizó un protocolo independiente de localización \\
\bottomrule
\end{tabular}
\end{table}

\section{Coherencia entre simulación y hardware}
\label{sec:comparacion_sim_real}

En la trayectoria rectilínea de 2~m, la desviación respecto de la referencia
nominal fue 3.39\% en simulación, calculada con la odometría, y 7.18\% en
hardware real, calculada con la distancia física medida
(Tabla~\ref{tab:resumen_odometria}). Por ello ambos valores no miden la misma
magnitud. El error odométrico solo se separa del error de seguimiento en
hardware real, donde alcanza 3.63\%, porque en simulación la pose verdadera de
Gazebo no quedó registrada.
En el protocolo angular, la desviación media respecto de la referencia
nominal de $360^\circ$ es mayor en simulación ($17.1^\circ$) que en hardware
real ($11.7^\circ$), con una dispersión menor ($5.8^\circ$ frente a
$8.0^\circ$), según la Tabla~\ref{tab:resumen_giro}.

El mapeo y la exploración se evaluaron únicamente sobre hardware real, por lo
que no admiten un contraste directo simulación--hardware. En hardware real, el
índice de cobertura libre promedió 62.2\% ($\sigma=12.9$) sobre las cinco
iteraciones válidas, con un área libre mapeada media de 3.50~m$^2$ y una mejor
iteración de 82.9\%. Las limitaciones metodológicas de cada ensayo se detallan
en el \ref{anx:resultados_metodologia}.
Por su parte, la caracterización cuantitativa del consumo energético y de la carga
computacional no forma parte de los ensayos realizados y se plantea como
trabajo futuro; su alcance se precisa en el \ref{anx:resultados_metodologia}.

\section*{Conclusión del capítulo}
\label{sec:conclusion_resultados}

Los resultados muestran un error odométrico medio de 3.63\% en
trayectorias rectas sobre hardware real y una desviación angular media de
$11.7^\circ$ respecto de la referencia nominal del protocolo. La respuesta en lazo cerrado de la Fig.~\ref{fig:respuesta_escalon_final} fue
medida en hardware, y sobre esa misma captura se calcularon los indicadores de
la Tabla~\ref{tab:pid_performance}, que muestran seguimiento de la referencia
sin sobreimpulso apreciable.
Por su parte, el sistema de mapeo con \texttt{slam\_toolbox} generó mapas utilizables en
hardware real. La exploración por fronteras alcanzó un índice
medio de cobertura libre de 62.2\% ($\sigma=12.9$) en cinco iteraciones, con
un área libre mapeada media de 3.50~m$^2$. La integración funcional y la
cobertura quedaron verificadas. La navegación hacia metas fijas definidas por
el operador queda fuera del alcance declarado.